\section{Zusammenfassung und Ausblick}
\label{sec:fazit}

Das Projekt beantwortet die Ausgangsfrage nicht mit einer einzelnen Kennzahl,
sondern mit einem nachvollziehbaren Analyseweg. Der gerichtete Netzwerkgraph
zeigt Deutschlands bilaterale Rolle zu einer Stunde. Die gestapelte Zeitreihe
verbindet absolute Erzeugung, relative Anteile, Preis und einen ausgewählten
Länderfluss. Die Pixelmatrix verdichtet 8184 physische Flusswerte zu einem
Überblick und dient zugleich als zweidimensionales Auswahlwerkzeug. Ein
gemeinsamer Elm-Zustand hält Land, Stunde und Zeitfenster konsistent.

Für den Mai 2025 zeigt die Fallanalyse deutliche Preisunterschiede und einen
starken negativen linearen Zusammenhang zwischen Preis und erneuerbarer
Leistung. Gleichzeitig ist der Zusammenhang zwischen Preis und einfacher Summe
der bilateralen physischen Flüsse schwach. Einzelne Länderbeziehungen können
zudem kommerziell und physisch verschiedene Richtungen aufweisen. Damit liefert
die Anwendung einen echten Mehrwert: Sie macht einfache Erklärungen prüfbar und
führt von einem Muster unmittelbar zu konkreten absoluten Werten.

Sinnvolle Erweiterungen sind:

\begin{itemize}
  \item frei wählbare Start- und Enddaten sowie serverseitig erzeugte Exporte,
  \item Drill-down von den vier Gruppen zu einzelnen Erzeugungstechnologien,
  \item ein ergänzender Scatterplot für Preis--Erzeugungs-Zusammenhänge,
  \item Vergleich physischer und kommerzieller Flüsse als Differenzansicht,
  \item saisonale oder mehrjährige Aggregation mit Zoom und
        Kalender-Hierarchie,
  \item eine Nutzerstudie mit Studierenden und energiewirtschaftlich
        Vorgebildeten.
\end{itemize}

Für längere Zeiträume müsste die Matrix aggregieren oder hierarchisch zoomen;
sonst werden Stunden schmaler als wahrnehmbare Pixel. Au{\ss}erdem sollte eine
künftige Analyse Nachfrage, installierte Leistung, Wetter und Netzengpässe
einbeziehen, bevor kausale Aussagen über Preis oder Fluss entstehen. In der
aktuellen Form bleibt das System bewusst ein exploratives Werkzeug mit klar
begrenzter Aussagekraft.
